\begin{problem}[第32届全国中学生物理竞赛复赛][圆环碰撞与斜抛运动]{42}
    如图，一质量分布均匀、半径为 $r$ 的刚性薄圆环落到粗糙的水平地面前的瞬间，圆环质心速度 $v_{0}$ 与竖直方向成 $\theta$（$\frac{\pi}{2} < \theta < \frac{3\pi}{2}$）角，并同时以角速度 $\omega_{0}$（$\omega_{0}$ 的正方向如图中箭头所示）绕通过其质心、且垂直环面的轴转动。已知圆环仅在其所在的竖直平面内运动，在弹起前刚好与地面无相对滑动，圆环与地面碰撞的恢复系数为 $k$，重力加速度大小为 $g$。忽略空气阻力。
    \begin{subproblem}
        求圆环与地面碰后圆环质心的速度和圆环转动的角速度；
    \end{subproblem}
    \begin{subproblem}
        求使圆环在与地面碰后能竖直弹起的条件和在此条件下圆环能上升的最大高度；
    \end{subproblem}
    \begin{subproblem}
        若让 $\theta$ 角可变，求圆环第二次落地点到首次落地点之间的水平距离 $s$ 随 $\theta$ 变化的函数关系式、$s$ 的最大值以及 $s$ 取最大值时 $r$、$v_{0}$ 和 $\omega_{0}$ 应满足的条件。
    \end{subproblem}
\end{problem}